\enteteExam{NFP 107 -- Systèmes de gestion de bases de données - Durée 2h30}%
{Session de rattrapage -- 2 septembre 2020}%
{Effectué à distance}%


Voici une partie de la base de données utilisée pour gérer des {\oe}uvres
dans un musée:
 
\begin{itemize}
  \item Oeuvre(\textbf{idOeuvre}, nom, idPropriétaire, idAuteur)
  \item Personne(\textbf{idPersonne}, nom, prénom)
  \item Expertise(\textbf{idExpertise}, idOeuvre, idExpert,  valeur)
\end{itemize}

Les auteurs, les propriétaires et les experts sont des personnes bien entendu. Les clés primaires sont en gras,
les clés primaires ne sont pas explicitement indiquées.


\section{Schéma relationnel (6 points)}

\begin{enumerate}

\item Pour chaque table, donner les dépendances fonctionnelles liant les clés primaires et clés étrangères (1 point)

\item En déduire les associations de type  ``plusieurs à 1'' du schéma Entité/Association correspondant à ce schéma 
    relationnel (1 point)

\item Donner ce schéma Entité/Association complet (1 point)


\item Donner des commandes SQL permettant de créer les tables
\texttt{\textbf{Oeuvre}} et \texttt{\textbf{Personne}}. (1 point) 
  
 \item J'ajoute la contrainte ``Un expert ne peut évaluer une {\oe}uvre qu'une 
    seule fois''. Quelle est la dépendance fonctionnelle correspondante
     et quel changement du schéma relationnel proposez-vous\,?  (1 point) 
     
    \item Ce schéma permet-il de représenter le fait qu'une personne peut être à la fois  propriétaire et auteur d'une
    même  {\oe}uvre\,? Argumentez votre réponse.
     (1 point) 


\end{enumerate}
\begin{solution}


\begin{itemize}
\item Dépendances entre clés primaires et clés étrangères :
   \begin{enumerate}
      \item Dans \texttt{Oeuvre}: $idOeuvre \to idAuteur$ et $idOeuvre \to idProprietaire$
       \item Dans \texttt{Expertise}: $idExpertise \to idOeuvre$ et $idExpertise \to idExpert$
  \end{enumerate}
   \item Chaque DF correspond à une association plusieurs à 1 (cf l'algorithme de normalisation). 
       On a donc ce type d'association entre Oeuvre et Personne, avec la sémantique "Auteur",
       entre Oeuvre et Personne une seconde fois mais avec la sémantique "Propriétaire",
         entre Expertise et Oeuvre et entre Expertise et Personne.
        \item Le schéma EA se déduit immédiatement de ce qui précède. Il n'y a pas d'association
        plusieurs à plusieurs.

 \item
  \begin{minted}{sql}
  CREATE TABLE Personne (
  idPersonne INTEGER NOT NULL,
  nom VARCHAR NOT NULL,
  prénom VARCHAR NOT NULL,
    PRIMARY KEY (idPersonne)
     );
   
  CREATE TABLE Oeuvre (
  idOeuvre INTEGER NOT NULL,
  nom VARCHAR,
  idAuteur INTEGER NOT NULL,
  idPropriétaire INTEGER NOT NULL,
   PRIMARY KEY (idOeuvre),
   FOREIGN KEY (idAuteur) REFERENCES Personne (idPersonne) 
   FOREIGN KEY (idPropriétaire) REFERENCES Personne (idPersonne) 
   );
  \end{minted}
  
  \item On a donc la nouvelle dépendance $idExpert \to idOeuvre$. Du coup la table Expertise
      n'est plus en 3FN. Un changement possible est que $idExpert$ devienne la clé
      de Expertise. 
      \item Oui, rien n'empêche  idAuteur et idPropriétaire d'être égaux dans la table Oeuvre
\end{itemize}
\end{solution}


\section{SQL (7 points)}

\begin{enumerate}
  \item Donnez le nom du propriétaire et le nom de l'auteur pour l'{\oe}uvre dont l'identifiant est 'fg65'.

  \item Donnez les noms des {\oe}uvres expertisées par  leur propriétaire
  
\item Donner les noms et prénoms des propriétaires d'une {\oe}uvre qui n'ont jamais effectué d'expertise.

\item Donnez les {\oe}uvres et le nom des personnes qui n'en sont ni propriétaire, ni auteur.

\item Donner les noms et prénoms des personnes qui sont à la fois auteur, propriétaire et expert (mais pas forcément
     de la même {\oe}uvre). 

\item Pour chaque {\oe}uvre donnez la moyenne des valeurs estimées par les experts.

\item Donnez le nombre d'expertises pour les {\oe}uvres dont la valeur maximale estimée est de 10\,000 Euros

  
\end{enumerate}

\begin{enumerate}
  \item \begin{minted}{sql}
     select p1.nom as nomAuteur, p2.nom as nomPropriétaire
     from Oeuvre as o, Personne as p1, Personne as p2
     where idOeuvre='fg65'
     and p.idAuteur=p1.idPersonne
     and p.idPropriétaire = p2.idPersonne
  \end{minted}
  
  \item 
  \begin{minted}{sql}
     select p1.nom as nomAuteur, p2.nom as nomPropriétaire
     from Oeuvre as o, Expertise as e
     where o.idOeuvre=e.idOeuvre
     and o.idPropriétaire = e.idPersonne
  \end{minted}
  
   \item \begin{minted}{sql}
     select  p.nom, p.prénom
     from Oeuvre as o, Personne as p
     where p.idPropriétaire = p.idPersonne
     and p.idPersonne not in (select idExpert from Expertise)
  \end{minted}
  
     \item \begin{minted}{sql}
     select  o.nom as nomOeuvre, p.nom as nomPersonne
     from Oeuvre as o, Personne as p
     where o.idPropriétaire != p.idPersonne
     and   o.idAuteur != p.idPersonne
  \end{minted}
  
     \item \begin{minted}{sql}
     select  distinct p.nom
     from  Personne as p, Oeuvre as o1, Oeuvre as o2, Expertise as e
     where o1.idPropriétaire = p.idPersonne
     and    o2.idAuteur = p.idPersonne
     and   e.idExpert = p.idPersonne 
  \end{minted}
  
       \item \begin{minted}{sql}
     select  o.nom, avg(valeur)
     from  Oeuvre as o, Expertise as e
     where o.idOeuvre = e.idOeuvre
     group by o.idOeuvre, o.nom
  \end{minted}
  
        \item \begin{minted}{sql}
     select  o.nom, count(*)
     from  Oeuvre as o, Expertise as e
     where o.idOeuvre = e.idOeuvre
     group by o.idOeuvre, o.nom
     having max(valeur) <= 10000
  \end{minted}
\end{enumerate}


\section{Un peu de logique (3 points)}

On dispose de deux prédicats $Auteur(x)$ et $Prop(x)$ qui sont vrais si,
respectivement, $x$ est auteur ou $x$ est propriétaire.

\begin{enumerate}
  \item Quelle formule exprime la condition ``Soit $x$  est propriétaire, soit $x$ est auteur mais pas les deux''

  \item Quelle est la négation de l'énoncé\,: ``Soit $x$ est propriétaire, soit $x$ est auteur''

  \item Comment exprimer l'énoncé suivant sans implication\,: ``Si $x$ est auteur, alors $x$ n'est pas propriétaire''
     (utiliser uniquement les connecteurs de SQL: and, or et not).
  
\end{enumerate}

\begin{solution}

\begin{enumerate}
  \item $Auteur(x) \lor Prop(x) \land \neg (Auteur(x) \land Prop(x))$

  \item $\neg Auteur(x) \land \neg Prop(x)$

  \item L'implication $p \to q$ est équivalente à $\neg p \lor q$.
       Donc en SQL on écrirait cette formule : $not(Auteur(x)) \lor Prop(x)$
  
\end{enumerate}


\end{solution}


\section{Algèbre (2 points)}
On dispose de deux tables $T_1(A,B,C)$ et $T_2(D,E,F)$. Donnez une  expression algébrique équivalente 
à la suivante, avec les contraintes suivantes\,: on peut utiliser  la jointure mais pas le  produit cartésien, et 
une sélection doit s'apppliquer directement à une table. 

$$\sigma_{A=C \land C > B \land E =F} (T_1 \times T_2)$$

\begin{solution}
$$\sigma_{A=C \land C > B} (T_1) \Join \sigma_{A=C \land E > F}(T_2)$$
\end{solution}

\section{Transactions (2 points)}

On dispose d'une table $T (id, valeur)$. Initialement toutes les valeurs sont différentes. 
Voici une procédure qui échange les valeurs de 2 nuplets.

\begin{minted}{sql}
 create or replace procedure Echange (id1 INT, id2 INT) AS
  -- Déclaration des variables
  val1, val2 integer;
  begin
    -- On recherche la valeur de id1 et de id2
    select valeur INTO val1 FROM T WHERE id = id1
    select valeur INTO val2 FROM T WHERE id = id2

    -- On échange les valeurs
    UPDATE T SET valeur = val1 WHERE id = id2
    UPDATE T SET valeur = val2 WHERE id = id1

 end;
\end{minted}

On est en mode \texttt{Autocommit}\,: un \texttt{commit} a lieu après chaque requête SQL. Expliquez
dans quel scénario l'exécution concurrente de deux procédures d'échange peut aboutir
à ce que deux nuplets aient la même valeur.


\begin{solution}
La première procédure s'exécute jusqu'à la fin du premier UPDATE. Un \texttt{commit} automatique
a lieu: les deux nuplets sont alors égaux. Si la seconde procédure commence alors, elle trouvera
une base avec deux nuplets égaux, ce qui ne devrait jamais arriver avec nos hypothèses en mode sérialisable.
\end{solution}

\end{document}

